\documentclass[10pt,a4paper]{amsart}
\usepackage[margin=19mm]{geometry}
\usepackage{amsmath,amssymb,mathtools}
\usepackage[scr=rsfs,cal=euler]{mathalfa}
\usepackage{xfrac}
\usepackage[unicode,hidelinks,bookmarks=false]{hyperref}
\newcommand{\ie}{i.e.\ }
\newcommand{\cf}{cf.\ }
\newcommand{\resp}{respectively\ }
\newcommand{\Subsection}{\subsection}
\providecommand{\nobreakdash}{\nobreak\hbox{-}}
\setlength{\emergencystretch}{2em}
\multlinegap0pt
\usepackage{mathtools}
\usepackage{amssymb}
\usepackage{tikz}
\usepackage[shortlabels]{enumitem}
\newtheorem{lemma}{Lemma}
\newcommand{\G}{\Gamma}
\newcommand{\Hom}{\operatorname{Hom}}
\newcommand{\X}{\mathcal X}
\newcommand{\id}{\mathrm{id}}
\newcommand{\modcat}{\ensuremath{\mbox{\rm -mod}}}
\newcommand{\modu}{\operatorname{mod}}
\title{Relative interval tilting, higher Auslander staircase corners and rational Dyck posets}
\author{Shengyong Pan}
\date{Source: arXiv:2608.06696v2 (CC BY 4.0)}
\begin{document}
\maketitle

\noindent\emph{Source and licence.} Relative interval tilting, higher Auslander staircase corners and rational Dyck posets. Version arXiv:2608.06696v2 (CC BY 4.0), distributed by arXiv under the Creative Commons Attribution 4.0 International licence, http://creativecommons.org/licenses/by/4.0/, as recorded in the arXiv licence metadata for this exact version and shown on the abstract page https://arxiv.org/abs/2608.06696v2. Full licence notice: \texttt{licenses/arxiv-2608.06696-CC-BY-4.0.txt}.\par\smallskip

\noindent\textit{Editorial scope.} Counted unit: the article's lemma lem:kan (source0.tex lines 1001--1016). The article's complete original proof of it is delivered with it (source0.tex lines 1017--1156, one \texttt{proof} environment of 140 lines). No mathematics is written, repaired or re-derived: the statement and the proof are the article's own lines, in the article's own notation, and no retained line is substituted or re-flowed.  No further source-local context is needed: every cross-reference of every retained line resolves inside the two delivered spans.  Nothing was omitted from any retained span, so the package carries no omission record.  No retained line contains a citation, so the wrapper carries no bibliography and no bibliography entry.  No retained line contains a figure environment, an image-inclusion command or drawing code, so no image is delivered and the package declares no asset dependency.  Editorial formatting only, in addition: an AMS \texttt{amsart} preamble that restates the article's own declarations and macros used by the retained lines, namely the environments \texttt{lemma} on separate counters, together with the standard packages those lines need.  The retained environments print in this document as Lemma 1 in order of appearance, whereas the article prints the same items with the numbering of its own full document.  The complete native source of the article as distributed in the arXiv e-print for this exact version is bundled unchanged as \texttt{sources/207025/source0.tex}.
\begin{lemma}
\label{lem:kan}
For each $y\in Y$, the functor $\iota_y^*$ has an exact right adjoint
\[
 (\iota_y)_*: F(y)\modcat\longrightarrow\G\modcat
\]
given on objects by
\begin{equation}
 ((\iota_y)_*M)(a,b)=
 \begin{cases}
 M(a),&b\leq y,\\
 0,&b\nleq y.
 \end{cases}
\label{eq:kan-formula}
\end{equation}
\end{lemma}

\begin{proof}
We first verify that the formula defines a $\G$-module. Let $(a,b)\in\operatorname{Ob}\G$.
If $b\leq y$, then $a\in F(b)\subseteq F(y)$
by the monotonicity of $F$. Hence $M(a)$ is defined.

Now let $h:(a,b)\longrightarrow(c,d)$ be nonzero only if
$b\leq d$, and its $\X$-component is a morphism
$ u:a\longrightarrow c$
in $\X$. Define
\[
((\iota_y)_*M)(h)
=
\begin{cases}
M(u), & d\leq y,\\
0,    & d\nleq y.
\end{cases}
\]
If $d\leq y$, then $b\leq d\leq y$, and therefore
\[
a\in F(b)\subseteq F(y),
\qquad
c\in F(d)\subseteq F(y).
\]
Since $F(y)$ is full in $\X$, the morphism $u:a\to c$ belongs to
$F(y)$, so $M(u)$ is defined.

The only apparently problematic case would be
\[
b\nleq y
\qquad\text{and}\qquad
d\leq y.
\]
This case cannot occur, because $b\leq d\leq y$ would imply
$b\leq y$. Thus the formula is well defined on every morphism.

It also respects identities. If $b\leq y$, the identity of $(a,b)$ is
sent to $M(\id_a)=\id_{M(a)}$.
If $b\nleq y$, the value at $(a,b)$ is the zero vector space, whose
identity morphism is the zero map.

To check composition, consider morphisms
\[
(a,b)\xrightarrow{h}(c,d)\xrightarrow{k}(e,r).
\]
If $r\leq y$, then $b\leq d\leq r\leq y$,
so all three objects lie in layers below $y$, and functoriality follows
from $M(kh)=M(k)M(h)$.
If $r\nleq y$, then the map associated with $kh$ is zero. The composite
of the two maps defined above is also zero, because the second map has
zero target. Therefore \eqref{eq:kan-formula} defines a covariant
$\G$-module.

We next prove the adjunction. Let $N\in\modu\G$ and
$M\in\modu F(y)$. Restriction to the layer $y$ defines a map
\[
\Phi:
\Hom_{\G}\bigl(N,(\iota_y)_*M\bigr)
\longrightarrow
\Hom_{F(y)}(\iota_y^*N,M).
\]
Explicitly, for a natural transformation
\[
\eta:N\longrightarrow(\iota_y)_*M,
\]
the component of $\Phi(\eta)$ at $a\in F(y)$ is
\[
\Phi(\eta)_a=\eta_{(a,y)}:
N(a,y)\longrightarrow M(a).
\]
Naturality of $\eta$ with respect to morphisms in the layer $y$ shows
that these components form a natural transformation $\iota_y^*N\longrightarrow M$.


Conversely, suppose that $\alpha:\iota_y^*N\longrightarrow M$
is a natural transformation. For every $(a,b)\in\operatorname{Ob}\G$
with $b\leq y$, let $v_{a,b}:(a,b)\longrightarrow(a,y)$
be the vertical morphism induced by $\id_a$. Define
\[
\Psi(\alpha)_{(a,b)}
=
\begin{cases}
\alpha_a\circ N(v_{a,b}), & b\leq y,\\
0,                         & b\nleq y.
\end{cases}
\]
We claim that these maps define a natural transformation
\[
\Psi(\alpha):N\longrightarrow(\iota_y)_*M.
\]

Indeed, let $h:(a,b)\longrightarrow(c,d)$
have $\X$-component $u:a\to c$. If $d\leq y$, then
$b\leq d\leq y$, and the following equality holds in $\G$:
\[
v_{c,d}\,h
=
\iota_y(u)\,v_{a,b}.
\]
Using this equality and the naturality of $\alpha$, we obtain
\[
\begin{split}
\Psi(\alpha)_{(c,d)}N(h)
&=
\alpha_cN(v_{c,d})N(h)\\
&=
\alpha_cN(\iota_y(u))N(v_{a,b})\\
&=
M(u)\alpha_aN(v_{a,b})\\
&=
((\iota_y)_*M)(h)\Psi(\alpha)_{(a,b)}.
\end{split}
\]
If $d\nleq y$, both sides of the required naturality identity are zero.
Thus $\Psi(\alpha)$ is natural.

Finally, $\Phi$ and $\Psi$ are inverse to each other. Since
$v_{a,y}=\id_{(a,y)}$, we have $\Phi(\Psi(\alpha))_a=\alpha_a$.

Conversely, if $\eta:N\to(\iota_y)_*M$ and $b\leq y$, naturality of
$\eta$ with respect to $v_{a,b}$ gives
\[
\eta_{(a,b)}
=
\eta_{(a,y)}N(v_{a,b}),
\]
because $(\iota_y)_*M$ sends $v_{a,b}$ to $\id_{M(a)}$. Hence $\Psi(\Phi(\eta))=\eta$.
We have therefore obtained a natural isomorphism
\[
\Hom_{\G}\bigl(N,(\iota_y)_*M\bigr)
\cong
\Hom_{F(y)}(\iota_y^*N,M),
\]
which proves that $(\iota_y)_*$ is right adjoint to $\iota_y^*$.

It remains to prove exactness. Kernels and cokernels in a functor
category are computed pointwise. By \eqref{eq:kan-formula}, evaluation
of $(\iota_y)_*$ at any object $(a,b)$ is either evaluation of $M$ at
$a$ or the zero functor. Both operations are exact. Therefore
$(\iota_y)_*$ is exact.
\end{proof}

\end{document}
